El análisis y diseño de algoritmos es un pilar fundamental en las Ciencias de la Computación, dándonos herramientas teóricas y prácticas para la creación de software eficiente y escalable. La relevancia de este campo radica en la resolución de problemas computacionales de manera óptima. En cuanto se habla del mundo académico se ha establecido un marco sólido para evaluar la eficiencia de los algoritmos a través de la notación de complejidad asintótica $\mathcal{O}$, que describe su comportamiento en el peor de los casos. Este nos permite clasificar y comparar algoritmos teóricamente, como es el caso de métodos de ordenamiento clásicos como \textbf{Merge Sort} y \textbf{Quick Sort} que poseen una complejidad promedio de $\mathcal{O}(n \log n)$, o la multiplicación de matrices con el algoritmo \textbf{Strassen}, que mejora la complejidad del método tradicional de $\mathcal{O}(n^3)$ a $\mathcal{O}(n^{2.81})$ \parencite{cormen2009introduction}.

\medskip

A pesar de que el análisis asintótico nos ofrece una sólida aproximación teórica, en realidad no captura el rendimiento real de una implementación en un escenario práctico. Factores como la arquitectura del procesador, el lenguaje de programación, la estructura de los datos de entrada (por ejemplo, si el arreglo está pre-ordenado o es aleatorio) y las constantes ocultas en la notación de complejidad pueden generar diferencias significativas en los tiempos de ejecución y el uso de memoria. Esto plantea una pregunta clave: ¿Cómo se comparan empíricamente algoritmos con complejidades teóricas similares o diferentes cuando se enfrentan a conjuntos de datos variados y de gran escala?

\medskip

El propósito de este informe es explorar esta brecha entre la teoría y la práctica mediante un análisis experimental de diferentes algoritmos de ordenamiento y de multiplicación de matrices. Para ello, se han implementado y medido los tiempos de ejecución y el uso de memoria de algoritmos como \textbf{Merge Sort} y \textbf{Quick Sort}, junto a la función \textit{std::sort} de C++. De esa misma forma, se compara el algoritmo \textbf{Naive} de multiplicación de matrices con la implementación del algoritmo de \textbf{Strassen}. Este trabajo además busca aplicar el uso de la metodología experimental (generación de diversos conjuntos de datos y la visualización de resultados), así permitiéndonos recoger información del comportamiento teórico predicho por la literatura con la evidencia empírica obtenida \parencite{musser1997introspective}.